\documentclass[11pt]{scrartcl}

\usepackage[stixtwo]{fontsetup}
\usepackage{amsmath}

\begin{document}
$x^2 - x - 6 = (x - 3)(x + 2)$, tehát
$D_f = \mathbb{R} \setminus \{-2, 3\}$. A függvény folytonos az
értelmezési tartományának minden pontjában, mert folytonos
függvények hányadosa, és a nevező nem nulla.
\[
	\lim_{x \to -2 \pm 0} f(x)
	= \lim_{x \to -2 \pm 0} \frac{\sin(x + 2)}{|x + 2|} \cdot \frac{1}{|x - 3|}
	=
	\lim_{x \to -2 \pm 0}
	\underbrace{\frac{\sin(x + 2)}{\pm (x + 2)}}_{\to \pm 1} \cdot
	\underbrace{\frac{1}{|x - 3|}}_{\to \frac{1}{5}}
	=
	\pm \frac{1}{5},
\]
tehát a függvénynek a $-2$ pontban véges ugrás típusú szakadása van.
\[
	\lim_{x \to 3} f(x)
	= \lim_{x \to 3}
	\underbrace{\frac{\sin(x + 2)}{|x + 2|}}_{\to \frac{\sin(5)}{5} < 0} \cdot
	\underbrace{\frac{1}{|x - 3|}}_{\to +\infty}
	= -\infty,
\]
azaz a $3$ pontban a bal és a jobb oldali határérték is $-\infty$, így itt a
függvénynek másodfajú szakadása van.
\end{document}
